\chapter*{ВВЕДЕНИЕ}
\addcontentsline{toc}{chapter}{ВВЕДЕНИЕ}

При использовании различных источников и устройств съемки получаются аэрофотоснимки с разным разрешением, также множество факторов могут привести к ухудшению их качества. Сведение этих аэрофотоснимков к единому разрешению позволяет унифицировать данные и комбинировать данные с разных источников, что существенно повышает точность и стабильность алгоритмов, работающих с обработкой и анализом визуальной информации.

\textbf{Цель выпускной квалификационной работы:} разработать метод сведения аэрофотоснимков с разными разрешениями к одному разрешению с использованием нейронной сети.

Для достижения поставленной цели требуется выполнить следующие задачи:

\begin{itemize}[label=---]
    \item провести анализ предметной области;
    \item рассмотреть существующие методы сведения аэрофотоснимков с разными разрешениями и сравнить их;
    \item разработать метод сведения аэрофотоснимков с разными разрешениями к одному разрешению с использованием нейронной сети;
    \item разработать программное обеспечение, реализующее описанный метод;
    \item исследовать зависимость качества изображения от коэффициента масштабирования.
\end{itemize}